\section{Introduction}

Leibnitz proposed that some time in the future, mathematics will be developed enough to model and solve every possible problem. Then came Russel

\begin{definition}
    \textbf{Russels Paradox:} Let $S$ be the set of sets that do not contain themselves.
    $$S = { X | X \notin X}$$
    Russels Paradox poses the following Question: Does $S$ contain intself. The answer is paradoxical
    $$S \notin S \implies S \in S$$
    $$S \in S \implies S \notin S$$
\end{definition}

\begin{remark}
    Hilbert style deduction allows you to solve Russels paradox by first stating what the axioms and interference rules of the System are.
\end{remark}

\begin{theorem}
    \textbf{Goedels incompletness:} There are mathematical statements that are true but can not possibly be proven to be so.
\end{theorem}

\begin{definition}
    \textbf{Turing-Machine (TM):} A theoretical machine consisting of an infinity strand of paper, a perfect pen and a set of rules. Every step includes reading current position, modifying it, moving left/right.
\end{definition}

\begin{remark}
    It can be shown that every algorithm can be modeld as a Turing-Machine and vice-versa.
\end{remark}

\begin{theorem}
    \textbf{Turing-Theorem:} There are problems that can not be automatically solved using an algorithm/TM.
\end{theorem}

\begin{proof}
    Let $w_{i,1}, w_{i,2}, w_{i,3} \dots$ be the output of Turing Machine $i$ for inputs $1,2,3 \dots$. 
    We define a new Turing Machine $M_x$ such that $w_{x,j} \neq w_{j,j}$ hence the output of $M_x$ differs from the output of every other Turing machine for at least one input. 
    Observe that this is a valid Turing Machine but it can not possibly be contained in our current list of Turing Machines. Hence, we reach a contradiction: if $M_x$ were in our enumeration at some index $x$, its output on input $x$ would be $w_{x,x} \neq w_{x,x}$, which is logically impossible. 
    Therefore, our assumption that we can construct such a machine $M_x$ must be false. This implies there is no computable general algorithm capable of determining the output $w_{j,j}$ for all Turing machines $j$, proving that there exist uncomputable functions and that the Halting Problem is undecidable.
\end{proof}

\begin{example}
    \textbf{Halting-problem:} Does algorithm $A$ terminate?
\end{example}


\begin{definition}
    \textbf{P vs. NP:} The p/np problem poses the question if all in polynomial time verifiable problems (NP) have a polynomial time solution (P).
\end{definition}
